% =====================================================================
% Figure 2 subsection -- REVISED 2026-09-08
% Numbers pulled from data/count_distribution/figure_2/<pair>/<mode>/comparison_summary.json
%   theta_hat / mean-variance   <- qc_filtered      (full gene panels, QC'd both sides)
%   total counts / genes-per-obs <- qc_and_hvg_matched (both sides HVG-subset to 556 genes)
% Configs from code/count_distribution/{cell,bin,bin16um,spot}_vs_*.sh
% =====================================================================

\subsection{Replication of real-world data}
Next, we evaluate whether the simulator can reproduce key statistical properties observed in real-world spatial transcriptomics data across measurement technologies and spatial resolutions. Our objective is to demonstrate the statistical flexibility of the simulator and its ability to be parameterized to approximate datasets generated using different technologies. We consider representative datasets generated using Xenium, Visium HD, and Visium. Spatial transcriptomics count data are typically overdispersed, with gene-level variance exceeding the mean for many genes, motivating the use of a negative binomial model for count generation \cite{zhao_modeling_2022} \cite{malsot_striping_2026} \cite{jones_cell_2025}. \\

Because gene panel sizes differ by orders of magnitude between the simulated marker panel (556 genes) and the real references (Xenium: 392 genes; Visium HD and probe-based Visium: $\sim$18{,}000 genes), each comparison is made on a matched gene set appropriate to the statistic. Gene-level mean--variance relationships are estimated on the full QC-filtered panel of each dataset, since highly-variable-gene selection mechanically biases the fitted dispersion; per-observation total-count and gene-detection distributions are instead compared after subsetting both datasets to a common 556-gene highly-variable set, so that differences in panel size do not dominate the comparison. Off-tissue observations are removed from both simulated and real data before comparison, mirroring the \texttt{in\_tissue} filtering applied to real Visium and Visium HD data. \\

We first compare simulated cell-resolution measurements with Xenium lung tissue datasets from non-diseased and lung cancer samples (\textbf{Figure~\ref{fig:figure_2}A}). Drawing baseline gene-expression values from a log-normal distribution with parameters $(-2.3, 0.7)$ and setting the negative-binomial dispersion to $\theta=0.40$ (with per-gene dispersion jitter $0.15$) yields an empirical dispersion estimate of $\hat{\theta}\approx 0.08$ for the simulated data, the same order as the estimates from the two Xenium datasets ($\hat{\theta}\approx 0.16$); on the mean--variance plot the fitted negative-binomial curves for simulated and real data are nearly coincident and both lie well above the Poisson expectation. Per matched observation, the simulated cells recover a comparable total count (median $153$ versus $73$--$107$ in the two Xenium samples) with a modestly higher zero fraction ($0.94$ versus $0.84$--$0.87$). The overall count scale can be tuned independently of dispersion through the baseline-expression and batch parameters, illustrating that dispersion and total molecular abundance are controlled separately within the simulation framework. \\

%% FIGURE 2 -------------------
\begin{figure}[!t]
\centering
\includegraphics[width=0.95\textwidth]{figure_2.png}
\caption{
\textbf{Comparison of simulated and real-world spatial transcriptomics data across measurement resolutions.}
(\textbf{A}) Cell-resolution simulations compared with Xenium lung datasets (non-diseased and lung cancer).
(\textbf{B}) 8~\textmu m bin-resolution simulations compared with a Visium HD breast cancer dataset.
(\textbf{C}) 16~\textmu m bin-resolution simulations compared with the same Visium HD breast cancer dataset resampled at 16~\textmu m.
(\textbf{D}) Spot-resolution simulations compared with a probe-based Visium (CytAssist FFPE) human lymph node dataset.
For each measurement resolution, the left panel compares the gene-level mean--variance relationship between simulated and real data (full QC-filtered gene panels), together with fitted negative-binomial relationships and the Poisson expectation ($\mathrm{variance}=\mathrm{mean}$). The middle panel compares the distribution of total molecular counts per observation, and the right panels compare the distributions of nonzero expression values before normalization, after library-size normalization to a target sum of 10{,}000 counts per observation, and after $\log(1+x)$ transformation; these per-observation comparisons are made after subsetting both datasets to a common 556-gene highly-variable set.
}
\label{fig:figure_2}
\end{figure}

We next evaluated the simulated data at Visium HD bin resolution using a breast cancer dataset, at both 8~\textmu m (\textbf{Figure~\ref{fig:figure_2}B}) and 16~\textmu m (\textbf{Figure~\ref{fig:figure_2}C}) bin sizes. To increase molecular density at bin resolution, we increased the three-dimensional cell-packing density by reducing the radius of the simulated tissue sphere ($R=2050$~\textmu m, $\sim$4\% packing) while holding the number of simulated cells fixed, rather than the sparser packing used for the cell-resolution simulation ($R=6000$~\textmu m); the same denser packing is used for the spot-resolution simulation below. This raised the expected number of molecules captured per spatial bin and produced mean--variance relationships in the negative-binomial regime at both bin sizes, with the simulated data somewhat less overdispersed than the real data ($\hat{\theta}\approx 0.28$ at 8~\textmu m and $\hat{\theta}\approx 0.18$ at 16~\textmu m, versus $\hat{\theta}\approx 1.2$ and $\hat{\theta}\approx 1.7$ for the corresponding real datasets). The 16~\textmu m comparison is the closest match on per-observation statistics: after gene-panel matching, the simulated and real bins have nearly identical median total counts ($176$ versus $179$), similar gene detection per bin (median $71$ versus $80$ genes), and similar zero fractions ($0.87$ versus $0.85$). At 8~\textmu m, panel matching is less informative because the 556 most-variable genes of an $\sim$18{,}000-gene real panel are only sparsely expressed in any single small bin, whereas the simulated 556-gene panel is the complete gene set: after matching, the simulated bins detect more per observation than the real bins ($485$ versus $35$ median counts, $140$ versus $25$ genes). \\

Finally, we compared simulated spot-resolution data (100~\textmu m spacing, 27.5~\textmu m capture radius) against a probe-based Visium human lymph node dataset (\textbf{Figure~\ref{fig:figure_2}D}). The simulated data reproduces the real mean--variance relationship well ($\hat{\theta}\approx 1.2$ versus $\hat{\theta}\approx 2.4$; fitted negative-binomial curves nearly coincident) and closely matches the real per-spot count statistics after gene-panel matching: median total counts $2355$ versus $2468$, median genes detected per spot $412$ versus $394$, and zero fraction $0.28$ versus $0.34$. A residual population of low-count simulated spots ($\sim$10--40 counts) persists from spots that fall on sparsely populated capture area; this is a geometric feature of the spot-placement grid and is not removed by adjusting the expression parameters. \\

Across measurement resolutions, the simulated data recapitulated the transition from approximately Poisson-like variability among low-abundance genes to increasing overdispersion among more highly expressed genes, and reproduced the real per-observation library-size and gene-detection distributions after accounting for the difference in gene panel size between the simulated marker panel and the real references. Together, these results demonstrate that the simulator can be tuned to approximate key distributional properties of spatial transcriptomics data across cell-, bin-, and spot-level measurement resolutions while retaining a fully known underlying spatial and molecular ground truth.
